\chapter{Implementation Program}

This chapter turns the preceding conceptual story into a concrete implementation program
for the repository as it exists now.

\section*{Principle 1: Keep the plain system as the stable substrate}
The foundational hierarchy in \repo{src/deppy/types/} should remain the canonical place for
symbolic type-level structure: substitution, free-variable logic, refinements, dependent
formers, indexed families, and subtyping. Hybrid features should continue to depend on
that language rather than re-implementing its semantics in parallel.

\section*{Principle 2: Concentrate trust logic in a small number of core modules}
The current repository exposes both lightweight and heavyweight trust vocabularies. This is
not fatal, but it suggests an engineering direction: preserve public convenience exports,
while concentrating the real semantics of trust, provenance, and evidence composition in a
small number of clearly canonical modules. That will make papers, dashboards, and checkers
less likely to drift apart.

\section*{Principle 3: Treat covers as first-class products}
The older monographs already emphasize cover synthesis and site families. The generation
agent and hybrid pipeline should continue to move in that direction. Rather than treating
module plans or verification stages as opaque steps, they should make covers, overlap
obligations, and gluing summaries first-class artifacts. This makes the local-to-global
story inspectable by users and not only by implementers.

\section*{Principle 4: Make residuals visible}
The paper has repeatedly insisted on epistemic honesty, and the implementation should do
the same. Residual obligations, assumption boundaries, and heuristic semantic judgments
should remain visible in manifests, renderers, and reports. The trust lattice is useful
precisely because it avoids silent collapse.

\section*{Principle 5: Let theory packs grow the scope}
The phrase ``absolutely everything'' should be realized by extensibility, not by claiming
that every domain is already solved. New theory packs, domain bridges, and site families
should remain the preferred path for broadening \DepPy{}'s scope. The repository already
has the architectural shape required for this growth.

\begin{remark}
The implementation program described here is conservative in one sense and radical in
another. It is conservative because it preserves the repository's existing semantic core.
It is radical because it proposes to explain more and more of the repository through a
single local-to-global language rather than through disconnected subsystem-specific
vocabularies.
\end{remark}
